%% 5-5 位置→速度→加速度（第11回：ページ制限なし。説明・図・例題を書き足した）
\bkchapter{5}{5}{位置→速度→加速度}{\bklevel{3}}{位置→速度→加速度}
\bkmichishirube{2-10,5-1,5-4}{5-13,6-1,6-15}

\begin{goal}
\begin{itemize}[label={\color{bkTheme}●}, leftmargin=1.6\zw, itemsep=0pt, topsep=1pt]
\item \bklevel{3}\enspace 速度と速さのちがい（向き）を説明できる。
\item \bklevel{3}\enspace 位置を2回微分して、加速度を求められる。
\item \bklevel{3}\enspace $\dot{x}$、$\ddot{x}$、$\dfrac{\dd^2x}{\dd t^2}$ の書き方を読める。
\end{itemize}
\end{goal}

\begin{hatena}
\noindent\begin{minipage}[c]{0.49\linewidth}
エレベーターに乗ると、上へ動き出すとき、体がぐっと重く感じます。止まるときは、ふわっと軽く感じます。でも、同じ速さで動いているあいだは、止まっているときと同じで、何も感じません。

体は、速さそのものではなく、何を感じているのでしょう。
\end{minipage}\hfill
\begin{minipage}[c]{0.48\linewidth}\centering
\begin{tikzpicture}[x=1mm, y=1mm]
  % 第11回（X14）：カゴは上面のまん中のロープで、機械室の滑車（巻き上げ機）から吊るす。滑車の反対側につり合いおもり
  % 昇降路のかべ
  \draw[bkGray, line width=0.9pt] (0,0) -- (0,47) (32,0) -- (32,47);
  \draw[bkGray, line width=0.6pt] (0,0) -- (32,0);
  % 機械室
  \draw[bkGray, line width=0.9pt, fill=bkGray!10] (0,47) rectangle (32,64);
  \fill[white] (10,46.4) rectangle (12,47.6) (24,46.4) rectangle (26,47.6);
  % 巻き上げ機（モーター）と滑車
  \draw[bkInk!70, line width=0.6pt, fill=bkGray!45, rounded corners=0.8] (1.6,51.5) rectangle (7,58.5);
  \draw[bkInk!70, line width=1pt] (7,55) -- (18,55);
  \draw[bkInk!80, line width=0.8pt, fill=bkGray!25] (18,55) circle (7);
  \fill[bkInk!80] (18,55) circle (1);
  % ロープ：カゴの上面のまん中 → 滑車の左 → 上を回って → 滑車の右 → おもり
  \draw[bkInk, line width=0.9pt] (11,29) -- (11,55) arc[start angle=180, end angle=0, radius=7] -- (25,35);
  % カゴ
  \draw[bkInk!80, line width=0.8pt, fill=bkPaper] (2,4) rectangle (20,29);
  \fill[bkInk!80] (9.8,29) rectangle (12.2,30.4);
  \bkperson{(9,6)}{14}{bkRoute6}
  \node[font=\sffamily\scriptsize, anchor=north east, text=bkGray, inner sep=1pt] at (19.2,28.2) {カゴ};
  % つり合いおもり
  \draw[bkInk!80, line width=0.7pt, fill=bkGray!55] (22.8,24) rectangle (27.2,35);
  % 名前
  \node[font=\sffamily\scriptsize, anchor=west, text=bkGray] at (0.8,61.6) {機械室};
  \node[font=\sffamily\scriptsize, anchor=north, text=bkGray, inner sep=1pt] at (25.4,22.6) {おもり};
  % 動く向きとラベル（かべの外）
  \draw[bk arrow, bkOrange] (36,6) -- (36,40);
  \node[anchor=west, font=\sffamily\scriptsize, align=left] at (38,40) {止まる\\軽い};
  \node[anchor=west, font=\sffamily\scriptsize, align=left] at (38,23) {一定の速さ\\ふつう};
  \node[anchor=west, font=\sffamily\scriptsize, align=left] at (38,6) {動き出す\\重い};
\end{tikzpicture}
\end{minipage}
\end{hatena}

体が感じているのは、速さそのものではなく、「速さが変わること」です。この章では、速さが変わるようすを数で表す方法を学びます。位置を微分すると速度、速度をもう一度微分すると加速度になります。

\clearpage
%% ---------------- やってみよう・思い出そう ----------------
\begin{yattemiyou}[車の速さの変わり方を調べよう]
止まっていた車が走り出し、しばらく同じ速さで走ってから、ブレーキをかけました。スピードメーターの値（km/h）を1秒ごとに記録しました。
\begin{enumerate}[label={\bkcircled[bkGreen]{\arabic*}}, leftmargin=2.4\zw, itemsep=1pt, topsep=3pt]
\item スピードメーターの値を、m/s になおそう（2-10）。たとえば\qty{36}{km/h}は、
\[ \qty{36}{km/h}\times\qty{1000}{m/km}\div\qty{3600}{s/h}=\qty{10}{m/s} \]
\qty{1000}{m/km} は「1 km あたり 1000 m」、\qty{3600}{s/h} は「1 時間あたり 3600 秒」という意味です。km どうし、h どうしが約分されて消え、m/s が残ります。
\item 1秒ごとに、速さがどれだけ変わったか（あとの速さ${}-{}$まえの速さ）を計算しよう。
\end{enumerate}
\begin{bkcenter}
\begin{bktable}{colspec={Q[l,m,wd=21mm] *{7}{Q[c,m,wd=7mm]} *{2}{Q[c,m,wd=8.6mm]}}, rowsep=1.5pt, colsep=1.1mm}
時刻 $t$（s） & 0 & 1 & 2 & 3 & 4 & 5 & 6 & 7 & 8 \\
メーター（km/h） & 0 & 9 & 18 & 27 & 36 & 36 & 36 & 25.2 & 14.4 \\
速さ $v$（m/s） & 0 & 2.5 & 5 & 7.5 & 10 & 10 & 10 & 7 & 4 \\
1秒ごとの変化 & & 2.5 & 2.5 & 2.5 & 2.5 & 0 & 0 & $-3$ & $-3$ \\
\end{bktable}
\bktabcaption{4行めは「あとの速さ${}-{}$まえの速さ」（単位は m/s）。}
\end{bkcenter}
\end{yattemiyou}

% 第7回：左の段の幅を字の幅の整数倍にして、字間がのびないようにする
\noindent\begin{minipage}[c]{14\zw}
走り出したときは、1秒ごとに速さが\qty{2.5}{m/s}ずつふえています。同じ速さのときは0、ブレーキをかけたときは、1秒ごとに\qty{3}{m/s}ずつへっています。速さと時間のグラフ（$v$-$t$ グラフ、2-3）にすると、この「1秒ごとの変化」は、グラフの傾きです。
\end{minipage}\hfill
\begin{minipage}[c]{0.57\linewidth}\centering
\begin{tikzpicture}
\begin{axis}[bkplot, width=70mm, height=40mm, xmin=0, xmax=8.6, ymin=0, ymax=12.5, xtick={1,...,8}, ytick={2.5,5,7.5,10}, xlabel={$t$（s）}, ylabel={$v$（m/s）}]
  \addplot[bkcurve, mark=*, mark size=1.4pt] coordinates {(0,0) (1,2.5) (2,5) (3,7.5) (4,10) (5,10) (6,10) (7,7) (8,4)};
  \node[font=\sffamily\scriptsize, text=bkOrange!80!black, anchor=north west] at (axis cs:2.3,4.6) {傾き $2.5$};
  \node[font=\sffamily\scriptsize, text=bkOrange!80!black, anchor=south] at (axis cs:5,10.7) {傾き $0$};
  \node[font=\sffamily\scriptsize, text=bkOrange!80!black, anchor=south west] at (axis cs:7.2,8) {傾き $-3$};
\end{axis}
\end{tikzpicture}
\bkfigcaption{やってみようの車の $v$-$t$ グラフ}
\end{minipage}

\begin{omoidasou}
\bkrecall{5-4}{$(x^n)'=nx^{n-1}$、定数を微分すると0。定数倍・和は1つずつ微分する。}
\bkrecall{5-1}{位置が小さくなる向きに動くとき、$\dfrac{\Delta x}{\Delta t}$ はマイナス。向きをふくめた速さが速度。}
\bkrecall{2-2}{マイナスは「反対の向き」を表す。$+$ の向きは、はじめに決めて変えない。}
\bkrecall{2-10}{換算の数も単位つきの量（\qty{1000}{m/km}、\qty{3600}{s/h}）。\\単位も計算できる。$(\mathrm{m/s})\div\mathrm{s}=\mathrm{m/s^2}$}
\bkrecall{2-3}{速さが一定なら、$v$-$t$ グラフの長方形の面積が道のり。}
\end{omoidasou}

%% ---------------- 数学の道具 ----------------
\bksection{doguu}{数学の道具}

\bksubsection{1}{速度は、位置を微分したもの}

\leavevmode\label{bk:5-5:sokudo}まっすぐな道を走る車の位置を、数直線の上の $x$（m）で表します。はじめに、右を $+$ の向きと決めます（2-2）。

\begin{bkfigure}
\begin{tikzpicture}[x=1mm, y=1mm]
  % 上：右へ動く車
  \draw[bkInk, line width=0.7pt, -{Stealth[length=3mm]}] (0,22) -- (122,22);
  \foreach \k/\l in {10/0,40/10,70/20,100/30} {\draw[bkInk, line width=0.6pt] (\k,20.8) -- (\k,23.2); \node[font=\sffamily\scriptsize, anchor=north] at (\k,20.3) {\l};}
  \node[font=\sffamily\scriptsize, anchor=north] at (119,20.3) {$x$（m）};
  \bkcar{(34,24.5)}{0.55}{bkBlue}
  \draw[bk arrow, bkOrange] (52,30) -- (68,30);
  \node[font=\sffamily\footnotesize, anchor=west, align=left] at (72,31) {右へ動く：$\Delta x>0$、速度は $+$};
  % 下：左へ動く車
  \draw[bkInk, line width=0.7pt, -{Stealth[length=3mm]}] (0,0) -- (122,0);
  \foreach \k/\l in {10/0,40/10,70/20,100/30} {\draw[bkInk, line width=0.6pt] (\k,-1.2) -- (\k,1.2); \node[font=\sffamily\scriptsize, anchor=north] at (\k,-1.7) {\l};}
  \node[font=\sffamily\scriptsize, anchor=north] at (119,-1.7) {$x$（m）};
  \begin{scope}[xscale=-1]\bkcar{(-88,2.5)}{0.55}{bkGreen}\end{scope}
  \draw[bk arrow, bkOrange] (70,8) -- (54,8);
  \node[font=\sffamily\footnotesize, anchor=east, align=right] at (50,9) {左へ動く：$\Delta x<0$、速度は $-$};
\end{tikzpicture}
\bkfigcaption{$+$ の向き（右）に動くと $\Delta x$ は正、反対の向きに動くと $\Delta x$ は負。}
\end{bkfigure}

$x$ がふえる向きに動けば、$\Delta x$ も $\dfrac{\Delta x}{\Delta t}$ も正です。$x$ がへる向きに動けば、$\Delta x$ も $\dfrac{\Delta x}{\Delta t}$ も負になります。この向きもふくめた瞬間の速さを\textbf{速度}といい、$v$ で表します。5-3で学んだように、瞬間の値は、位置 $x$ を時間 $t$ で微分すると求まります。\bkidx{そくど}{速度}
\[ v=\frac{\dd x}{\dd t} \]
速度から向き（$+$ か $-$ か）を取った大きさが、ふだん「速さ」とよんでいるものです。

\begin{bkcenter}
\begin{bktable}{colspec={Q[c,m,wd=26mm] Q[c,m,wd=26mm] Q[c,m,wd=26mm]}, rowsep=1.5pt}
速度 & 向き & 速さ \\
\qty{10}{m/s} & $+$ の向き（右） & \qty{10}{m/s} \\
$\qty{-10}{m/s}$ & $-$ の向き（左） & \qty{10}{m/s} \\
\qty{0}{m/s} & 止まっている & \qty{0}{m/s} \\
\end{bktable}
\bktabcaption{速度が $\qty{10}{m/s}$ と $\qty{-10}{m/s}$ は、速さは同じで向きが反対。}
\end{bkcenter}

\noindent\textbf{例：真上に投げたボール}\par
\noindent ボールを真上に、秒速\qty{20}{m}で投げ上げます。上向きを $+$ にすると、$t$ 秒後の高さは
\[ x=20t-5t^2\quad\text{（m）} \]
です。なぜこの式になるのでしょう。もし重力がなければ、ボールは秒速\qty{20}{m}のまま上がり続けて、$t$ 秒で $20t$（m）上がります。でも、実際には重力で引っぱられ、手をはなした物が落ちるのと同じだけ下がります。5-2のように、$t$ 秒で落ちる距離は $5t^2$（m）でした（ほんとうは $4.9t^2$ ですが、ここでも計算しやすいように5にします）。だから、高さは $20t$ から $5t^2$ を引いた $20t-5t^2$ になります。

微分すると（5-4）、
\[ v=\frac{\dd x}{\dd t}=20-5\times2t=20-10t \]

\noindent\begin{minipage}[c]{0.58\linewidth}
\begin{bktable}{colspec={Q[c,m,wd=17mm] *{5}{Q[c,m,wd=8.6mm]}}, rowsep=1.5pt, colsep=1.2mm, cells={font=\small}}
$t$（s） & 0 & 1 & 2 & 3 & 4 \\
$x$（m） & 0 & 15 & 20 & 15 & 0 \\
$v$（m/s） & 20 & 10 & 0 & $-10$ & $-20$ \\
\end{bktable}

\vspace{4pt}
1秒後の速度は $20-10=10$ で、上向きに\qty{10}{m/s}。3秒後は $20-30=-10$ で、下向きに\qty{10}{m/s}。2秒後は $v=0$ で、ボールは一瞬止まり、いちばん高い所（\qty{20}{m}）にいます。1秒後と3秒後は、高さも速さも同じですが、向きが反対です。速度を使うと、この向きのちがいが $+$ と $-$ で表せます。
\end{minipage}\hfill
\begin{minipage}[c]{0.38\linewidth}\centering
\begin{tikzpicture}[x=1mm, y=1.6mm]
  \draw[bkInk, line width=0.6pt] (0,0) -- (0,22);
  \foreach \h in {0,5,10,15,20} {\draw[bkInk, line width=0.5pt] (-1,\h) -- (1,\h); \node[font=\sffamily\scriptsize, anchor=east] at (-1.5,\h) {\h};}
  \node[font=\sffamily\scriptsize, anchor=south] at (0,22.5) {$x$（m）};
  % 上がるとき（左の列）と、下りるとき（右の列）
  \foreach \t/\h/\v in {0/0/20,1/15/10} {\shade[ball color=bkOrange] (8,\h) circle[radius=1.6mm]; \draw[bk arrow, bkBlue] (8,{\h+1.2}) -- (8,{\h+1.2+\v/3});}
  \shade[ball color=bkOrange] (14,20) circle[radius=1.6mm];
  \foreach \t/\h/\v in {3/15/10,4/0/20} {\shade[ball color=bkOrange] (20,\h) circle[radius=1.6mm]; \draw[bk arrow, bkRed] (20,{\h-1.2}) -- (20,{\h-1.2-\v/3});}
  \node[font=\sffamily\scriptsize, anchor=west] at (9.8,0) {0秒};
  \node[font=\sffamily\scriptsize, anchor=west] at (9.8,15) {1秒};
  \node[font=\sffamily\scriptsize, anchor=south] at (14,21.2) {2秒};
  \node[font=\sffamily\scriptsize, anchor=west] at (21.8,15) {3秒};
  \node[font=\sffamily\scriptsize, anchor=west] at (21.8,1.2) {4秒};
\end{tikzpicture}
\end{minipage}

\vspace{2pt}
\noindent{\small 右の図：青の矢印は上向きの速度、赤の矢印は下向きの速度。矢印が長いほど速い。}

\bksubsection{2}{加速度は、速度を微分したもの}

\leavevmode\label{bk:5-5:kasokudo}やってみようの車では、走り出したときに、1秒ごとに速度が\qty{2.5}{m/s}ずつふえていました。このように、速度が変わる速さを\textbf{加速度}（かそくど）といい、$a$ で表します。\bkidx{かそくど}{加速度}

\noindent\textbf{平均の加速度}\par
\noindent 時間 $\Delta t$ のあいだに速度が $\Delta v$ 変わったとき、$\dfrac{\Delta v}{\Delta t}$ を\textbf{平均の加速度}といいます。平均の速さ $\dfrac{\Delta x}{\Delta t}$（5-1）の $x$ を $v$ にかえただけです。やってみようの車の0秒から4秒までなら、
\[ \frac{\Delta v}{\Delta t}=\frac{\qty{10}{m/s}-\qty{0}{m/s}}{\qty{4}{s}}=\frac{\qty{10}{m/s}}{\qty{4}{s}}=\qty{2.5}{m/s^2} \]

\noindent\textbf{単位 \unit{m/s^2} の意味}\par
\noindent 速度の単位 m/s を、時間の単位 s でわるので、$(\mathrm{m/s})\div\mathrm{s}=\mathrm{m/s^2}$ です（2-10）。「メートル毎秒毎秒」と読みます。\qty{2.5}{m/s^2} は、「\textbf{1秒あたり、速度が\qty{2.5}{m/s}ずつ変わる}」という意味です。s が2回出てくるのは、1回めが速度の「毎秒」、2回めが「1秒ごとに変わる」の「毎秒」だからです。

\noindent\textbf{瞬間の加速度}\par
\noindent 平均の速さで $\Delta t$ を0に近づけると、瞬間の速度になりました（5-2・5-3）。同じように、平均の加速度で $\Delta t$ を0に近づけると、その瞬間の加速度になります。つまり、加速度は速度を時間で微分したものです。
\[ a=\frac{\dd v}{\dd t} \]
グラフで見ると、$x$-$t$ グラフの接線の傾きが速度 $v$、$v$-$t$ グラフの接線の傾きが加速度 $a$ です。やってみようの車の加速度は、走り出したときが \qty{2.5}{m/s^2}、一定の速さのときが \qty{0}{m/s^2}、ブレーキのときが $\qty{-3}{m/s^2}$ です。

\noindent\textbf{加速度のマイナス}\par
\noindent ブレーキのときの $\qty{-3}{m/s^2}$ は、速度が1秒ごとに\qty{3}{m/s}ずつへることを表しています。$+$ の向きに走っている車の加速度が $-$ なら、車はおそくなります（まちがえやすいところも見よう）。

\noindent\textbf{身近な乗り物の加速度}\par
\noindent 0から時速\qty{100}{km}まで10秒で加速する車を考えます。まず、速度を m/s になおします。
\[ \qty{100}{km/h}\times\qty{1000}{m/km}\div\qty{3600}{s/h}\approx\qty{27.8}{m/s} \]
（$100000\div3600=27.77\cdots$ を、小数第1位までのがい数にした。$\approx$ は「ほぼ等しい」を表す記号で、5-8 でくわしく学びます）
平均の加速度は、$\dfrac{\qty{27.8}{m/s}}{\qty{10}{s}}=\qty{2.78}{m/s^2}$ です。

\bksubsection{3}{はてな？の答え：体は加速度を感じている}

\leavevmode\label{bk:5-5:hatena}\bkhatenaref のエレベーターの速度を、グラフにしてみます。上向きを $+$ とします。このエレベーターは、2秒かけて\qty{2}{m/s}まで速くなり、6秒間その速さで上がり、2秒かけて止まります。

\begin{bkfigure}
\begin{tikzpicture}
\begin{axis}[bkplot, name=e1, width=96mm, height=34mm, xmin=0, xmax=10.8, ymin=0, ymax=2.9, xtick={2,4,6,8,10}, ytick={1,2}, xlabel={$t$（s）}, ylabel={$v$（m/s）}]
  \addplot[bkcurve] coordinates {(0,0) (2,2) (8,2) (10,0)};
  \node[font=\sffamily\scriptsize, text=bkOrange!80!black, anchor=north west] at (axis cs:1.15,0.95) {傾き1};
  \node[font=\sffamily\scriptsize, text=bkOrange!80!black, anchor=south] at (axis cs:5,2.1) {傾き0};
  \node[font=\sffamily\scriptsize, text=bkOrange!80!black, anchor=north east] at (axis cs:8.85,0.95) {傾き$-1$};
\end{axis}
\begin{axis}[bkplot, at={($(e1.south)+(0,-11mm)$)}, anchor=north, width=96mm, height=34mm, xmin=0, xmax=10.8, ymin=-1.7, ymax=1.7, xtick={2,4,6,8,10}, xticklabels={}, ytick={-1,1}, xlabel={$t$（s）}, ylabel={$a$（\unit{m/s^2}）}]
  \addplot[bkcurve3, line width=1.4pt] coordinates {(0,1) (2,1)};
  \addplot[bkcurve3, line width=1.4pt] coordinates {(2,0) (8,0)};
  \addplot[bkcurve3, line width=1.4pt] coordinates {(8,-1) (10,-1)};
  \node[font=\sffamily\scriptsize, anchor=west, inner sep=1pt] at (axis cs:2.3,1) {重く感じる};
  \node[font=\sffamily\scriptsize, anchor=south, inner sep=3pt] at (axis cs:5,0) {ふつう};
  \node[font=\sffamily\scriptsize, anchor=east, inner sep=1pt] at (axis cs:7.7,-1) {軽く感じる};
\end{axis}
\end{tikzpicture}
\bkfigcaption{上：エレベーターの $v$-$t$ グラフ。下：その傾き（加速度）のグラフ。}
\end{bkfigure}

速度のグラフ（上）は、0秒から8秒まで、ずっと $+$（上向き）です。でも、体の感じ方は3つに分かれます。感じ方が分かれる所は、下の加速度のグラフが分かれる所とぴったり同じです。
\begin{itemize}[leftmargin=1.2\zw, itemsep=1pt, topsep=2pt, label={\color{bkTheme}●}]
\item 動き出すとき（0〜2秒）：加速度は $+\qty{1}{m/s^2}$（上向き）。重く感じる。
\item 一定の速さのとき（2〜8秒）：加速度は0。止まっているときと同じ。
\item 止まるとき（8〜10秒）：加速度は $\qty{-1}{m/s^2}$（下向き）。軽く感じる。
\end{itemize}
体が感じていたのは、速度ではなく\textbf{加速度}だったのです。

\begin{motto}[3]
\textbf{体重計に乗って、エレベーターに乗ると}\par
エレベーターの中で体重計に乗ると、上へ動き出すときは目もりが大きく、止まるときは小さくなります。上向きに加速するには、床が足を、ふだんより強くおし上げなければならないからです。床から強くおされるのを、体は「重い」と感じます。くわしくは、力と運動を学ぶ第6巻で考えます。
\end{motto}

\bksubsection{4}{2回微分する}

\leavevmode\label{bk:5-5:nikai}加速度は、位置 $x$ を微分して速度 $v$ にし、その $v$ をもう一度微分したものです。これを、位置を\textbf{2回微分する}といいます。

\begin{bkfigure}
\begin{tikzpicture}[x=1mm, y=1mm, q/.style={bk node, minimum width=30mm, minimum height=12mm}]
  \node[q] (x) at (0,0) {位置 $x$\\（m）};
  \node[q] (v) at (50,0) {速度 $v$\\（m/s）};
  \node[q] (a) at (100,0) {加速度 $a$\\（\unit{m/s^2}）};
  \draw[bk link] (x) -- node[bk link label, above] {時間で微分} (v);
  \draw[bk link] (v) -- node[bk link label, above] {時間で微分} (a);
\end{tikzpicture}
\bkfigcaption{位置を1回微分すると速度、2回微分すると加速度。単位は、1回微分するごとに s でわった形になる。}
\end{bkfigure}

式で書くと、次のようになります。
\[ a=\frac{\dd v}{\dd t}=\frac{\dd}{\dd t}\Bigl(\frac{\dd x}{\dd t}\Bigr)=\frac{\dd^2x}{\dd t^2} \]
$\dfrac{\dd}{\dd t}$ は「時間で微分する」という命令の記号です。命令を2回つづけたので、上には $\dd$ が2つ、下には $\dd t$ が2つあります。これを短く書いたのが $\dfrac{\dd^2x}{\dd t^2}$ です。上の $\dd^2$ は「$\dd$ を2回」、下の $\dd t^2$ は「$\dd t$ を2回」という意味で、$t^2$ を微分するという意味ではありません。

\begin{bkcenter}
\begin{bktable}{colspec={Q[l,m,wd=26mm] Q[c,m,wd=18mm] Q[c,m,wd=18mm] Q[l,m,wd=58mm]}, rowsep=1.5pt, cells={font=\small}, row{1}={font=\sffamily\bfseries\small}}
書き方 & 1回微分 & 2回微分 & 読み方・使う分野 \\
数学 & $x'$ & $x''$ & ダッシュ、ツーダッシュ \\
ライプニッツ & $\dfrac{\dd x}{\dd t}$ & $\dfrac{\dd^2x}{\dd t^2}$ & ディーツー・エックス・ディーティーにじょう \\
ニュートン & $\dot{x}$ & $\ddot{x}$ & エックス・ドット、ツードット（物理で、時間の微分に使う） \\
\end{bktable}
\bktabcaption{1回微分は速度、2回微分は加速度。専門書では場面で使い分ける。}
\end{bkcenter}

ニュートンの書き方では、文字の上の点の数が、時間で微分した回数です。$\dot{x}$ は1回、$\ddot{x}$ は2回です。短く書けるので、物理の本でよく使われます。

\bkkigocards{\bkkigocard{$\dot{x}$、$\ddot{x}$}{エックス・ドット、エックス・ツードット}{時間で1回微分（速度）・2回微分（加速度）}{x-dot}{5-5}}
{\bkkigocard{$\dfrac{\dd^2x}{\dd t^2}$}{ディーツー・エックス・ディーティーにじょう}{$x$ を $t$ で2回微分したもの}{second derivative}{5-5}}

\noindent\textbf{例}\par
\noindent 位置が $x=t^3$（m）のとき、1回微分して $v=\dot{x}=3t^2$、もう1回微分して $a=\ddot{x}=3\times2t=6t$ です。2秒では、$v=3\times2^2=\qty{12}{m/s}$、$a=6\times2=\qty{12}{m/s^2}$ です。

\bksubsection{5}{落ちる物の加速度}

\leavevmode\label{bk:5-5:g}手をはなしたボールの、落ちた距離はほんとうは $x=4.9t^2$（m）です（下向きを $+$）。2回微分すると、
\[ v=\dot{x}=4.9\times2t=9.8t,\qquad a=\ddot{x}=9.8 \]
表で見ると、速度は1秒ごとに\qty{9.8}{m/s}ずつふえています。

\begin{bkcenter}
\begin{bktable}{colspec={Q[l,m,wd=30mm] *{4}{Q[c,m,wd=12mm]}}, rowsep=1.5pt}
$t$（s） & 0 & 1 & 2 & 3 \\
$x=4.9t^2$（m） & 0 & 4.9 & 19.6 & 44.1 \\
$v=9.8t$（m/s） & 0 & 9.8 & 19.6 & 29.4 \\
1秒ごとの $v$ の変化 & & 9.8 & 9.8 & 9.8 \\
\end{bktable}
\bktabcaption{落ちる物の速度は、1秒ごとに\qty{9.8}{m/s}ずつふえる。}
\end{bkcenter}

加速度は時間によらず、いつも \qty{9.8}{m/s^2} です。これを\textbf{重力加速度}といい、$g$ で表します。9.8という数は、計算で決めたものではなく、地球の上で、物を落としてはかった値です。月の上では、重力が弱いので、約\qty{1.6}{m/s^2}です。\bkidx{じゅうりょくかそくど}{重力加速度 $g$}

空気の抵抗がなければ、重い鉄の球も軽いプラスチックの球も、同じ加速度 $g$ で落ちます。空気のない月の上で、宇宙飛行士がハンマーと鳥の羽を同時に落とすと、2つは同時に地面に着きました（1971年、アポロ15号）。

真上に投げ上げたボール $x=v_0t-4.9t^2$（上向きを $+$。$v_0$ は投げたときの速度）でも、2回微分すると $a=-9.8$ です。上がっているときも、いちばん高い所でも、落ちているときも、加速度はいつも下向きに $g$ です。

\bksubsection{6}{グラフで見る：$x$・$v$・$a$}

\leavevmode\label{bk:5-5:graph}\noindent\begin{minipage}[c]{0.56\linewidth}
\begin{tikzpicture}
% 第6回：段の間を9mmにして、v・a の名前が上のグラフにかからないようにする。
% v のグラフの目もり数字は、線のない側に置く（1は下、2は左下、3・4は上）
\begin{axis}[bkplot, name=g1, width=66mm, height=28mm, xmin=0, xmax=4.4, ymin=0, ymax=24, xtick={1,2,3,4}, ytick={10,20}, xlabel={$t$}, ylabel={$x$}]
  \addplot[bkcurve, domain=0:4] {20*x-5*x^2};
\end{axis}
\begin{axis}[bkplot, at={($(g1.south)+(0,-9mm)$)}, anchor=north, name=g2, width=66mm, height=28mm, xmin=0, xmax=4.4, ymin=-22, ymax=24, xtick={1,2,3,4}, xticklabels={}, ytick={-20,20}, xlabel={$t$}, ylabel={$v$}]
  \addplot[bkcurve2, domain=0:4] {20-10*x};
  \node[anchor=north, font=\sffamily\footnotesize, text=bkInk, inner sep=2pt] at (axis cs:1,0) {$1$};
  \node[anchor=north east, font=\sffamily\footnotesize, text=bkInk, inner sep=1.5pt] at (axis cs:2,0) {$2$};
  \node[anchor=south, font=\sffamily\footnotesize, text=bkInk, inner sep=2.5pt] at (axis cs:3,0) {$3$};
  \node[anchor=south, font=\sffamily\footnotesize, text=bkInk, inner sep=2.5pt] at (axis cs:4,0) {$4$};
\end{axis}
\begin{axis}[bkplot, at={($(g2.south)+(0,-9mm)$)}, anchor=north, width=66mm, height=28mm, xmin=0, xmax=4.4, ymin=-14, ymax=6, xtick={1,2,3,4}, ytick={-10}, xlabel={$t$}, ylabel={$a$}]
  \addplot[bkcurve3, domain=0:4] {-10};
\end{axis}
\end{tikzpicture}
\end{minipage}\hfill
\begin{minipage}[c]{0.4\linewidth}
1で使った、投げ上げたボールの位置・速度・加速度のグラフです。
\[ x=20t-5t^2 \]
\begin{itemize}[leftmargin=1.2\zw, itemsep=1pt, topsep=2pt, label={\color{bkTheme}●}]
\item $x$-$t$ の接線の傾きが $v$。
\item $v$-$t$ の傾きが $a$。
\item 最も高い所（$t=2$）で $v=0$。でも $a$ は0ではなく、$-10$。
\end{itemize}
\end{minipage}

\vspace{2pt}
上のグラフ（$x$）は、山の形です。山をのぼっている0〜2秒は傾きが正なので、まん中のグラフ（$v$）は $+$ です。山のてっぺんの2秒で傾きは0、山を下りる2〜4秒は傾きが負なので、$v$ は $-$ です。まん中の $v$ のグラフは、どこでも同じ傾き $-10$ の直線です。だから、下の $a$ のグラフは、いつも $-10$ の横の直線になります。

\bksubsection{7}{加速度が一定の運動の式}

\leavevmode\label{bk:5-5:toukasoku}落ちるボールの式 $x=4.9t^2$ の4.9は、どこから来たのでしょう。また、加速度がいつも同じ値 $a$ の運動は、どんな式で表せるでしょう。ここでは、「2回微分すると $a$ になる式」を、微分を逆にたどって作ってみます。

\noindent\textbf{①\enspace 2回微分すると $a$ になる式}\par
\noindent $t^2$ を2回微分すると、$t^2\to2t\to2$ で、2になります。2ではなく $a$ にしたいので、$t^2$ に $\dfrac{a}{2}$ をかけます。
\[ \frac{1}{2}at^2\ \xrightarrow{\text{微分}}\ \frac{1}{2}a\times2t=at\ \xrightarrow{\text{微分}}\ a \]
$\dfrac{1}{2}$ は、$t^2$ を微分したときに前に出てくる2を、打ち消すためについています。

\noindent\textbf{②\enspace はじめの速度 $v_0$ を入れる}\par
\noindent ①の式の速度は $at$ なので、$t=0$ で速度が0です。はじめに速度 $v_0$ で動いていたときは、$v_0t$ をたします。$v_0t$ は、1回微分すると $v_0$、2回微分すると0なので、加速度は $a$ のまま変わりません。

\noindent\textbf{③\enspace はじめの位置 $x_0$ を入れる}\par
\noindent 定数 $x_0$ をたしても、微分すると0なので、速度も加速度も変わりません。$t=0$ での位置が $x_0$ になります。

まとめると、加速度が一定の値 $a$ の運動は、
\[ x=x_0+v_0t+\frac{1}{2}at^2,\qquad v=v_0+at \]
と表せます（この式の別の求め方は5-13）。$t=0$ を入れると、$x=x_0$、$v=v_0$ で、たしかに「はじめの位置」と「はじめの速度」になっています。

落ちるボールは、$x_0=0$、$v_0=0$、$a=9.8$ の場合です。$\dfrac{1}{2}\times9.8=4.9$ なので、$x=4.9t^2$ になります。\textbf{4.9は、重力加速度9.8の半分}だったのです。

\noindent\textbf{例：ブレーキで止まる車}\par
\noindent 秒速\qty{20}{m}で走る車が、加速度 $\qty{-4}{m/s^2}$ でブレーキをかけます。$x_0=0$、$v_0=20$、$a=-4$ なので、
\[ v=20-4t,\qquad x=20t+\frac{1}{2}\times(-4)\times t^2=20t-2t^2 \]

\begin{bkcenter}
\begin{bktable}{colspec={Q[c,m,wd=18mm] *{6}{Q[c,m,wd=10mm]}}, rowsep=1.5pt, cells={font=\small}}
$t$（s） & 0 & 1 & 2 & 3 & 4 & 5 \\
$v$（m/s） & 20 & 16 & 12 & 8 & 4 & 0 \\
$x$（m） & 0 & 18 & 32 & 42 & 48 & 50 \\
\end{bktable}
\bktabcaption{ブレーキをかけた車の速度と位置。速度は1秒ごとに\qty{4}{m/s}ずつへる。}
\end{bkcenter}

$v=0$ になるのは、$20-4t=0$ より $4t=20$、$t=5$ 秒後です。そのあいだに進む距離は、$x=20\times5-2\times5^2=100-50=\qty{50}{m}$ です。

\vspace{4pt}
\noindent\begin{minipage}[c]{0.5\linewidth}
\textbf{グラフでも確かめる}\par
速さが一定なら、$v$-$t$ グラフの長方形の面積が道のりでした（2-3）。速度がへっていくときも、グラフの下の面積が道のりになります（くわしくは5-9・5-13）。右の三角形の面積は、
\[ \frac{1}{2}\times\qty{5}{s}\times\qty{20}{m/s}=\qty{50}{m} \]
で、式の答えと同じです。
\end{minipage}\hfill
\begin{minipage}[c]{0.46\linewidth}\centering
\begin{tikzpicture}
\begin{axis}[bkplot, width=58mm, height=38mm, xmin=0, xmax=5.8, ymin=0, ymax=24, xtick={1,...,5}, ytick={10,20}, xlabel={$t$（s）}, ylabel={$v$（m/s）}]
  \fill[bkYellow!45] (axis cs:0,0) -- (axis cs:0,20) -- (axis cs:5,0) -- cycle;
  \addplot[bkcurve, domain=0:5] {20-4*x};
  \node[font=\sffamily\scriptsize, anchor=center] at (axis cs:1.4,5.5) {面積50};
\end{axis}
\end{tikzpicture}
\end{minipage}

\begin{motto}[3]
\textbf{3回微分すると}\par
加速度をもう一度微分した $\dddot{x}$ を、加加速度（ジャーク）といいます。加速度がどれだけ急に変わるかを表し、乗り物の乗り心地を設計するときに使います。はてな？のエレベーターのグラフでは、加速度が2秒と8秒でいきなり変わっています。実際のエレベーターは、加速度を少しずつ変えて、がくんとしないようにしています。
\end{motto}

\begin{taisetsu}
速度 $v=\dfrac{\dd x}{\dd t}=\dot{x}$（m/s）\qquad 加速度 $a=\dfrac{\dd v}{\dd t}=\dfrac{\dd^2x}{\dd t^2}=\ddot{x}$（\unit{m/s^2}）\\
位置を1回微分すると速度、2回微分すると加速度\\
加速度が一定なら\quad $x=x_0+v_0t+\dfrac{1}{2}at^2$、\enspace$v=v_0+at$
\end{taisetsu}

\begin{kotoba}
\small
\bkword{速度}{向きもふくめた速さ。位置を時間で微分したもの。}
\bkword{加速度}{速度の変わる速さ。速度を時間で微分したもの。単位は \unit{m/s^2}。}
\bkword{重力加速度}{落ちる物の加速度。地球の上で $g=\qty{9.8}{m/s^2}$。}
\end{kotoba}

%% ---------------- 例題 ----------------
\bksection{reidai}{例題でたしかめよう}

\begin{reidai}[速度と加速度を求める]
位置が $x=3t^2+2t$（m）の物体の、速度 $v$ と加速度 $a$ を求めましょう。\par
\kangaekata $x$ を $t$ で1回微分すると $v$、その $v$ をもう1回微分すると $a$ です。\par
\shiki $v=\dfrac{\dd x}{\dd t}=3\times2t+2\times1=6t+2$\par
\hspace*{2.6\zw}$a=\dfrac{\dd v}{\dd t}=6\times1+0=6$\par
\kotae $v=6t+2$（m/s）、$a=6$（\unit{m/s^2}）
\end{reidai}

\begin{jibun}
位置が $x=2t^3$（m）の物体の、$v$ と $a$ を求めましょう。また、1秒での $v$ と $a$ の値を答えましょう。
\bkjans{$v=6t^2$、$a=12t$。1秒で $v=\qty{6}{m/s}$、$a=\qty{12}{m/s^2}$}
\end{jibun}

\begin{reidai}[平均の加速度]
止まっていた車が、8秒で時速\qty{72}{km}になりました。平均の加速度を求めましょう。\par
\kangaekata 加速度の単位は \unit{m/s^2} なので、速度を m/s になおしてから、$\dfrac{\Delta v}{\Delta t}$ を計算します。\par
\shiki $\qty{72}{km/h}\times\qty{1000}{m/km}\div\qty{3600}{s/h}=\qty{20}{m/s}$\par
\hspace*{2.6\zw}$\Delta v=\qty{20}{m/s}-\qty{0}{m/s}=\qty{20}{m/s}$、$\Delta t=\qty{8}{s}$\par
\hspace*{2.6\zw}$\dfrac{\Delta v}{\Delta t}=\dfrac{\qty{20}{m/s}}{\qty{8}{s}}=\qty{2.5}{m/s^2}$\par
\kotae \qty{2.5}{m/s^2}
\end{reidai}

\begin{jibun}
止まっていた電車が、30秒で時速\qty{54}{km}になりました。平均の加速度を求めましょう。
\bkjans{$\qty{54}{km/h}\times\qty{1000}{m/km}\div\qty{3600}{s/h}=\qty{15}{m/s}$。$15\div30=\qty{0.5}{m/s^2}$}
\end{jibun}

\begin{reidai}[ある時刻の値]
位置が $x=t^3-6t^2+9t$（m）の物体の、2秒での速度と加速度を求めましょう。\par
\kangaekata まず $v$ と $a$ を $t$ の式で求めてから、$t=2$ を入れます。\par
\shiki $v=3t^2-6\times2t+9=3t^2-12t+9$\par
\hspace*{2.6\zw}$a=3\times2t-12=6t-12$\par
\hspace*{2.6\zw}$t=2$ で $v=3\times2^2-12\times2+9=12-24+9=-3$\par
\hspace*{2.6\zw}$a=6\times2-12=0$\par
\kotae 速度 $\qty{-3}{m/s}$、加速度 \qty{0}{m/s^2}
\end{reidai}

\begin{jibun}
位置が $x=t^3-3t$（m）の物体の、2秒での速度と加速度を求めましょう。
\bkjans{$v=3t^2-3$、$a=6t$。速度 \qty{9}{m/s}、加速度 \qty{12}{m/s^2}}
\end{jibun}

\begin{reidai}[投げ上げ]
真上に投げたボールの高さが $x=15t-5t^2$（m、上向きが $+$）のとき、（1）1秒と2秒の速度、（2）速度が0になる時刻、（3）加速度を求めましょう。\par
\kangaekata 速度の符号で向きがわかります。$+$ なら上向き、$-$ なら下向きです。\par
\shiki $v=15-5\times2t=15-10t$、\enspace$a=-10$\par
\hspace*{2.6\zw}（1）$t=1$ で $v=15-10=5$、$t=2$ で $v=15-20=-5$\par
\hspace*{2.6\zw}（2）$15-10t=0$ より $10t=15$、$t=1.5$\par
\hspace*{2.6\zw}（3）$a=-10$（時刻によらない）\par
\kotae （1）上向きに\qty{5}{m/s}、下向きに\qty{5}{m/s}\quad（2）1.5秒（最も高い所）\quad（3）$\qty{-10}{m/s^2}$
\end{reidai}

\begin{jibun}
高さが $x=25t-5t^2$（m）のボールについて、1秒と4秒の速度、速度が0になる時刻、加速度を求めましょう。
\bkjans{\qty{15}{m/s}、$\qty{-15}{m/s}$。2.5秒。$\qty{-10}{m/s^2}$}
\end{jibun}

\begin{reidai}[加速度が一定の運動の式をつくる]
はじめの位置が\qty{2}{m}、はじめの速度が\qty{3}{m/s}、加速度が \qty{4}{m/s^2} で一定の物体があります。位置 $x$ と速度 $v$ を $t$ の式で表し、2秒での位置と速度を求めましょう。\par
\kangaekata $x=x_0+v_0t+\dfrac{1}{2}at^2$、$v=v_0+at$ に、$x_0=2$、$v_0=3$、$a=4$ を入れます。\par
\shiki $x=2+3t+\dfrac{1}{2}\times4\times t^2=2+3t+2t^2$\par
\hspace*{2.6\zw}$v=3+4t$\par
\hspace*{2.6\zw}$t=2$ で $x=2+6+8=16$、$v=3+8=11$\par
\hspace*{2.6\zw}確かめ：$x$ を微分すると $3+4t$ で、$v$ と同じ\par
\kotae $x=2t^2+3t+2$、$v=4t+3$。2秒で\qty{16}{m}、\qty{11}{m/s}
\end{reidai}

\begin{jibun}
止まっていた物体が、はじめの位置0から、加速度 \qty{6}{m/s^2} で一定のまま動きはじめました。$x$ と $v$ を $t$ の式で表し、3秒での位置と速度を求めましょう。
\bkjans{$x=3t^2$、$v=6t$。3秒で\qty{27}{m}、\qty{18}{m/s}}
\end{jibun}

\begin{reidai}[ブレーキで止まるまで]
秒速\qty{15}{m}で走る自転車が、加速度 $\qty{-3}{m/s^2}$ でブレーキをかけました。止まるまでの時間と、そのあいだに進む距離を求めましょう。\par
\kangaekata 加速度が一定なので、$x=v_0t+\dfrac{1}{2}at^2$、$v=v_0+at$ を使います（$x_0=0$）。止まるのは $v=0$ のときです。\par
\shiki $v=15+(-3)\times t=15-3t$。$15-3t=0$ より $t=5$\par
\hspace*{2.6\zw}$x=15t+\dfrac{1}{2}\times(-3)\times t^2=15t-1.5t^2$\par
\hspace*{2.6\zw}$t=5$ を入れて $x=15\times5-1.5\times25=75-37.5=37.5$\par
\kotae 5秒、\qty{37.5}{m}
\end{reidai}

\begin{jibun}
秒速\qty{20}{m}で走る車が、加速度 $\qty{-5}{m/s^2}$ でブレーキをかけました。止まるまでの時間と距離を求めましょう。
\bkjans{$20-5t=0$ より4秒。$20\times4-2.5\times16=\qty{40}{m}$}
\end{jibun}

\noindent\textbf{速さが2倍になると、止まるまでの距離は？}\par
\noindent 加速度 $\qty{-5}{m/s^2}$ でブレーキをかけたとき、はじめの速さと、止まるまでの距離をくらべます（じぶんでとこう6と同じ計算）。
\begin{bkcenter}
\begin{bktable}{colspec={Q[l,m,wd=40mm] *{3}{Q[c,m,wd=16mm]}}, rowsep=1.5pt}
はじめの速さ（m/s） & 10 & 20 & 30 \\
止まるまでの時間（s） & 2 & 4 & 6 \\
止まるまでの距離（m） & 10 & 40 & 90 \\
\end{bktable}
\bktabcaption{速さが2倍になると距離は4倍、3倍になると距離は9倍。}
\end{bkcenter}
速さが2倍になると、止まるまでの時間が2倍になり、そのあいだの平均の速さも2倍になります。だから、距離は $2\times2=4$ 倍です。車のスピードを出しすぎると危ないのは、このためです。

\begin{reidai}[速度が0になる時刻が2つあるとき]
位置が $x=t^3-6t^2+9t$（m）の物体の速度が0になる時刻を求め、それぞれの時刻の加速度を答えましょう。\par
\kangaekata $v=0$ は $t$ の2次方程式になります。因数分解で解きます（3-8・3-10）。\par
\shiki $v=3t^2-12t+9=3(t^2-4t+3)=3(t-1)(t-3)$\par
\hspace*{2.6\zw}$v=0$ より $t-1=0$ または $t-3=0$。$t=1,\ 3$\par
\hspace*{2.6\zw}$a=6t-12$ より、$t=1$ で $a=6-12=-6$、$t=3$ で $a=18-12=6$\par
\kotae 1秒（加速度 $\qty{-6}{m/s^2}$）と3秒（加速度 \qty{6}{m/s^2}）
\end{reidai}

\begin{jibun}
位置が $x=t^3-3t^2$（m）の物体の速度が0になる時刻と、そのときの加速度を求めましょう。
\bkjans{$v=3t^2-6t=3t(t-2)$ より0秒と2秒。$a=6t-6$ より $\qty{-6}{m/s^2}$ と \qty{6}{m/s^2}}
\end{jibun}

\begin{reidai}[落ちる物の速さ]
落ちる物の位置は $x=4.9t^2$（m、下向きが $+$）です。速さが\qty{19.6}{m/s}になるのは何秒後ですか。そのとき、何m落ちていますか。\par
\kangaekata 速さ $v$ を $t$ の式で表し、$v=19.6$ となる $t$ を求めます。\par
\shiki $v=4.9\times2t=9.8t$。$9.8t=19.6$ より $t=19.6\div9.8=2$\par
\hspace*{2.6\zw}$x=4.9\times2^2=4.9\times4=19.6$\par
\kotae 2秒後、\qty{19.6}{m}（加速度はいつも $g=\qty{9.8}{m/s^2}$）
\end{reidai}

\begin{jibun}
落ちる物（$x=4.9t^2$）の速さが\qty{29.4}{m/s}になるのは何秒後ですか。そのとき、何m落ちていますか。
\bkjans{$9.8t=29.4$ より3秒後。$4.9\times9=\qty{44.1}{m}$}
\end{jibun}

%% ---------------- まちがえやすいところ・物理工学 ----------------
\bksection{tsumazuki}{まちがえやすいところをチェック}

\begin{tsumazuki}{速度が0なら、加速度も0？}
投げ上げたボールは、最も高い所で一瞬止まります（$v=0$）。でも、加速度は $-10$ のままです。もし加速度も0なら、速度は0のまま変わらないので、ボールは空中に止まったままになるはずです。速度が0でも、速度が変わり続けていれば、加速度は0ではありません。
\end{tsumazuki}

\begin{tsumazuki}{加速度がマイナスなら、おそくなる？}
加速度の符号は、速度が変わる「向き」です。速度も加速度もマイナスなら、マイナスの向きにどんどん速くなります。投げ上げたボールが落ちてくるとき（$v<0$、$a<0$）は、速くなっていきます。速くなるかおそくなるかは、$v$ と $a$ の符号が同じ（速くなる）か、ちがう（おそくなる）かで決まります。
\end{tsumazuki}

\begin{tsumazuki}{時速を秒速になおすとき、かける・わるを逆にする}
\bkng{$72\times3.6=259.2$ で、\qty{259.2}{m/s}}\bkok{$\qty{72}{km/h}\times\qty{1000}{m/km}\div\qty{3600}{s/h}=\qty{20}{m/s}$}
単位をつけて計算すると、km と h が約分されて m/s が残るかどうかで、まちがいに気づけます。秒速\qty{259}{m}は、新幹線よりずっと速く、車としておかしいことにも気づけます。
\end{tsumazuki}

\begin{butsuri}
\begin{itemize}[leftmargin=1.2\zw, itemsep=2pt, topsep=0pt, label={\color{bkPurple}●}]
\item \textbf{加速度センサー}\\
スマートフォンには、加速度をはかる小さな部品が入っています。画面の向きが変わるのは、重力が下向きに引っぱる向きを、この部品が感じとっているからです。車のエアバッグは、ぶつかったときの大きな加速度を感じて開きます。
\item \textbf{力と加速度}\\
物にはたらく力の合計は、質量×加速度に等しい（$F=ma$）。これがニュートンの運動の法則です。加速度は位置の2回微分なので、運動の法則は「2回微分をふくむ式」になります。これが第6巻の微分方程式です（6-15）。
\item \textbf{乗り心地}\\
加速度がいきなり変わると、乗っている人はがくんと感じます。エレベーターや電車は、加速度をゆっくり変えるように動かしています。
\end{itemize}
\end{butsuri}

%% ---------------- 章末問題 ----------------
\bkendproblems
\bkproblemgroup{基本}{3}

\bkproblem 位置が $x=4t^2-3t+2$（m）の物体の、$v$ と $a$ を求めましょう。
\bkpans{$v=8t-3$（m/s）、$a=8$（\unit{m/s^2}）}

\bkproblem 位置が $x=2t^3-3t^2$（m）の物体の、2秒での速度と加速度を求めましょう。
\bkpans{$v=6t^2-6t$ より \qty{12}{m/s}、$a=12t-6$ より \qty{18}{m/s^2}}

\bkproblem 真上に投げたボールの高さが $x=30t-5t^2$（m）のとき、（1）1秒と4秒の速度、（2）最も高くなる時刻、（3）加速度を求めましょう。
\bkpans{（1）\nobreak \qty{20}{m/s}、$\qty{-10}{m/s}$\quad（2）\nobreak 3秒\quad（3）\nobreak $\qty{-10}{m/s^2}$}

\bkproblem 止まっていた自転車が、5秒で時速\qty{36}{km}になりました。平均の加速度を求めましょう。
\bkpans{$\qty{36}{km/h}\times\qty{1000}{m/km}\div\qty{3600}{s/h}=\qty{10}{m/s}$、$10\div5=\qty{2}{m/s^2}$}

\bkproblem 時速\qty{72}{km}で走る車が、加速度 $\qty{-5}{m/s^2}$ でブレーキをかけました。止まるまでの時間と距離を求めましょう。
\bkpans{\qty{20}{m/s} なので、$20-5t=0$ より4秒、$20\times4-2.5\times16=\qty{40}{m}$}

\bkproblem 位置が $x=t^3-12t$（m）の物体について、速度が0になる時刻（$t>0$）と、そのときの加速度を求めましょう。
\bkpans{$v=3t^2-12=3(t-2)(t+2)$ より2秒。$a=6t$ より \qty{12}{m/s^2}}

\bkproblem 落ちる物（$x=4.9t^2$）の、3秒後の速さと加速度を求めましょう。
\bkpans{$v=9.8t$ より \qty{29.4}{m/s}。加速度 \qty{9.8}{m/s^2}}

\bkproblemgroup{チャレンジ}{3}

\bkproblem 位置が $x=t^3-9t^2+24t$（m）の物体について、速度が0になる時刻を2つ求め、それぞれの時刻の加速度を答えましょう（2次方程式は3-10）。
\bkpans{$v=3t^2-18t+24=3(t-2)(t-4)$ より2秒と4秒。$a=6t-18$ より、2秒で $\qty{-6}{m/s^2}$、4秒で \qty{6}{m/s^2}}

\bkproblem 空気の抵抗がないとき、落ちる物は $x=4.9t^2$（m）です。速さが時速\qty{100}{km}になるのは、落ちはじめてから何秒後ですか。小数第1位まで求めましょう。
\bkpans{時速\qty{100}{km}は約\qty{27.8}{m/s}。$9.8t=27.8$ より、約2.8秒後}

\bkproblem 止まっていた電車が、加速度 \qty{0.8}{m/s^2} で動きはじめました。時速\qty{72}{km}になるまでの時間と、そのあいだに進む距離を求めましょう。
\bkpans{\qty{20}{m/s}。$v=0.8t=20$ より25秒。$x=\dfrac{1}{2}\times0.8\times25^2=\qty{250}{m}$}

\bkproblem エレベーターが、止まった状態から3秒かけて上向きに\qty{1.5}{m/s}になり、7秒間その速さで上がり、2秒かけて止まりました。（1）それぞれのあいだの加速度を求めましょう。（2）体が重く感じるのは、どのあいだですか（はてな？の答え）。
\bkpans{（1）\nobreak 0〜3秒：\qty{0.5}{m/s^2}、3〜10秒：\qty{0}{m/s^2}、10〜12秒：$\qty{-0.75}{m/s^2}$\quad（2）\nobreak 0〜3秒（動き出すとき）}

\begin{furikaeri}
\begin{checklist}
\item 速度と速さのちがいを説明できる。\hfill{\footnotesize\sffamily\color{bkTheme}→\pageref{bk:5-5:sokudo}ページ}
\item 位置を2回微分して、加速度を求められる。\hfill{\footnotesize\sffamily\color{bkTheme}→\pageref{bk:5-5:nikai}ページ}
\item $\dot{x}$、$\ddot{x}$、$\dfrac{\dd^2x}{\dd t^2}$ を読める。\hfill{\footnotesize\sffamily\color{bkTheme}→\pageref{bk:5-5:nikai}ページ}
\end{checklist}
\end{furikaeri}

\bkfukidashi{sora}{$x^3$ や $x^4$ は微分できるようになったけど、$(2x+1)^5$ みたいな式は、展開しないとだめなのかな？}
